% !TeX program = pdflatex
\documentclass[11pt,a4paper]{article}
\usepackage{../../recursos/latex/estilo-notas}

\begin{document}
\cabecalhoaula{Extra E3}{NLP $+$ Classificação (Aplicação)}{Aplicação (notebook); AME §8.1.3 e Apêndice A.2}

%==============================================================================
\section*{Objetivos da aula}
%==============================================================================
Ao final desta aula o aluno deve ser capaz de:
\begin{itemize}[leftmargin=1.4cm]
  \item descrever as etapas de \textbf{limpeza} e exploração de uma base textual;
  \item transformar texto em atributos numéricos com \textbf{\emph{bag-of-words}} e
        \textbf{TF--IDF};
  \item escolher classificadores adequados a dados de texto (Bayes ingênuo,
        logística penalizada) e justificar a escolha;
  \item avaliar o classificador com as métricas certas em um problema
        \textbf{desbalanceado} (\emph{spam});
  \item montar o \textbf{\emph{pipeline}} completo, integrando o que foi visto no
        curso inteiro.
\end{itemize}

\begin{emsala}
Aula de \textbf{síntese}: um problema real (filtrar \emph{spam} de SMS) que costura
quase todas as aulas --- vetorização de texto, alta dimensão e esparsidade
(Aula~05), Bayes ingênuo (Aula~08), métricas em dados desbalanceados (Aula~09) e
\emph{pipelines} (Aula~07). Pergunta de abertura: \emph{``um modelo só entende
números. Como transformar uma mensagem de texto em algo que ele consiga
classificar?''}
\end{emsala}

%==============================================================================
\section{O problema: filtrar spam}
%==============================================================================
Usaremos a base \texttt{spam.csv} (SMS Spam Collection): $\sim$5\,574 mensagens de
SMS rotuladas como \emph{ham} (legítima, $Y=0$) ou \emph{spam} ($Y=1$). É um problema
de \textbf{classificação binária} (Aula~08) cujas covariáveis são \emph{texto bruto}
--- e o desafio central é a \emph{representação}: máquinas só processam números.

\begin{atencao}
A base é \textbf{desbalanceada} (muito mais \emph{ham} que \emph{spam}). Como vimos
na Aula~09, a acurácia engana aqui --- precisaremos de precisão, \emph{recall} e
$F_1$. E os custos são assimétricos: bloquear uma mensagem legítima (falso positivo)
costuma ser pior que deixar passar um \emph{spam}.
\end{atencao}

%==============================================================================
\section{Do texto aos atributos}
%==============================================================================
\subsection{Limpeza e normalização}
Texto bruto é ruidoso. As etapas usuais de pré-processamento:
\begin{itemize}[leftmargin=1.4cm,nosep]
  \item \emph{minusculizar} (\texttt{lower}): ``Free'' e ``free'' viram o mesmo
        \emph{token};
  \item remover \textbf{pontuação} e \textbf{números} (conforme o problema);
  \item \emph{tokenizar}: quebrar a mensagem em palavras (\emph{tokens});
  \item remover \textbf{\emph{stopwords}}: palavras muito frequentes e pouco
        informativas (``the'', ``is'', ``and'');
  \item \emph{stemming}/\emph{lematização} (opcional): reduzir flexões a um radical
        (``running'' $\to$ ``run'').
\end{itemize}

\subsection{\emph{Bag-of-words} e a matriz documento--termo}
A representação mais simples é o \textbf{\emph{bag-of-words}} (saco de palavras):
ignora-se a ordem e conta-se a ocorrência de cada palavra. Constrói-se a
\textbf{matriz documento--termo} $\bm{X}$, em que a linha $i$ é a mensagem $i$, a
coluna $j$ é um termo do vocabulário, e $x_{ij}$ é a contagem (ou presença binária)
do termo $j$ na mensagem $i$.

\begin{ideia}
Cada mensagem vira um vetor num espaço de dimensão $=$ tamanho do vocabulário ---
facilmente \emph{milhares} de covariáveis. Mas cada mensagem usa poucas palavras,
então $\bm{X}$ é \textbf{esparsa} (quase toda zero). Guarde-a em formato esparso
(Apêndice A.3 do [AME]) para economizar memória.
\end{ideia}

\begin{atencao}[Eis a Aula~05 em ação]
Texto é o exemplo canônico de \textbf{alta dimensão $+$ esparsidade}: milhares de
palavras, mas a maioria irrelevante para a classe. Por isso métodos que
\emph{selecionam variáveis} (Lasso, Bayes ingênuo) vão bem, e o KNN vai mal (não
seleciona; a distância é dominada por palavras irrelevantes) --- exatamente o que a
teoria da Aula~05 previu.
\end{atencao}

\subsection{TF--IDF}
Contagens cruas dão peso demais a palavras comuns. O \textbf{TF--IDF}
(\emph{term frequency--inverse document frequency}) reescala cada termo:
\[
  \text{tf-idf}(t,p) = \underbrace{\text{tf}(t,p)}_{\substack{\text{frequência do termo}\\\text{no documento }p}}
  \times
  \underbrace{\log\frac{N}{\text{df}(t)}}_{\substack{\text{raridade do termo}\\\text{na coleção}}},
\]
onde $\text{df}(t)$ é o número de documentos que contêm $t$ e $N$ o total. Palavras
que aparecem em \emph{todo} documento têm IDF baixo (pouco discriminantes); palavras
\emph{raras e específicas} ganham peso. Costuma melhorar a classificação de texto.

%==============================================================================
\section{Análise exploratória}
%==============================================================================
Antes de modelar, vale \emph{olhar} os dados: quais as palavras mais frequentes em
\emph{ham} vs.\ \emph{spam}? Em \emph{spam} surgem termos como ``free'', ``win'',
``call'', ``txt'', ``prize''; em \emph{ham}, vocabulário do dia a dia. Esse contraste
já sugere que o \emph{bag-of-words} carrega sinal suficiente para separar as classes.

%==============================================================================
\section{Classificando}
%==============================================================================
Quais métodos usar? Os que prosperam em alta dimensão esparsa:
\begin{description}[leftmargin=2.8cm,style=nextline]
  \item[Bayes ingênuo (multinomial)] o ``cavalo de batalha'' do texto. A suposição
        de independência condicional (Aula~08) é grosseira para palavras, mas
        funciona surpreendentemente bem e é rápida. A versão \emph{multinomial}
        modela contagens de palavras. (Lembre da Aula~08: trabalhe em escala
        logarítmica para evitar \emph{underflow}.)
  \item[Logística penalizada (Lasso/Ridge)] discriminativa, robusta, e o Lasso faz
        \emph{seleção de palavras} --- gerando um modelo interpretável (quais
        palavras pesam para \emph{spam}). É o que o [AME] usa no exemplo das resenhas
        da Amazon.
  \item[Florestas / \emph{boosting}] também aplicáveis; XGBoost costuma ir bem.
\end{description}

\begin{ideia}
Por que Bayes ingênuo combina tão bem com texto? Porque o modelo multinomial sobre
contagens de palavras é natural, há pouquíssimos parâmetros por palavra (baixa
variância) e a independência condicional, embora falsa, raramente troca a
\emph{classe} de maior probabilidade --- e, para classificar (Aula~09), basta acertar
o lado do corte, não a probabilidade exata.
\end{ideia}

%==============================================================================
\section{Avaliação}
%==============================================================================
Como o problema é desbalanceado (Aula~09):
\begin{itemize}[leftmargin=1.4cm,nosep]
  \item reporte a \textbf{matriz de confusão}, \textbf{precisão}, \textbf{\emph{recall}}
        e \textbf{$F_1$} --- não só a acurácia;
  \item para \emph{spam}, a \emph{precisão} sobre a classe ``spam'' é crítica
        (não classificar \emph{ham} como \emph{spam});
  \item ajuste o \textbf{corte} de decisão conforme os custos (Aula~09): se um falso
        positivo é muito caro, exija probabilidade alta para marcar como \emph{spam}.
\end{itemize}

\begin{emsala}[fechar o curso com uma métrica que discorda das outras duas]
A §8 da \texttt{Aula prática E3} acrescenta ao saco de palavras o comprimento da
mensagem e a proporção de dígitos --- duas colunas que separam muito ($70{,}9$
contra $138{,}8$ caracteres; $0{,}4\%$ contra $11{,}7\%$ de dígitos). Pergunte à
turma se o modelo melhora, e depois mostre: acurácia $0{,}9821\to0{,}9862$, AUC
$0{,}9879\to0{,}9950$ e AP $0{,}9691\to\mathbf{0{,}9572}$.

A AUC sobe porque as duas colunas arrumam o grosso do ranking; a AP olha a cabeça
dele, onde o TF--IDF já ia bem. Como o filtro bloqueia as mensagens de maior escore,
quem manda é a AP --- e por ela as colunas não pagam o que custam. É um bom
encerramento: a métrica tem de ser escolhida \emph{antes}, pela pergunta, ou os
mesmos três números justificam três decisões diferentes.
\end{emsala}

%==============================================================================
\section{O \emph{pipeline} completo (síntese do curso)}
%==============================================================================
Tudo se encaixa num único \emph{pipeline} (Aula~07): vetorização $\to$ classificador,
com hiperparâmetros escolhidos por validação cruzada (Aula~03) \emph{sem vazamento}
--- e o vetorizador deve aprender o vocabulário \textbf{só no treino}.

\begin{lstlisting}
import pandas as pd
from sklearn.feature_extraction.text import CountVectorizer, TfidfVectorizer
from sklearn.naive_bayes import MultinomialNB
from sklearn.linear_model import LogisticRegression
from sklearn.pipeline import Pipeline
from sklearn.model_selection import train_test_split, GridSearchCV
from sklearn.metrics import classification_report, confusion_matrix

sms = pd.read_csv("spam.csv", encoding="ISO-8859-1")
X, y = sms["text"], (sms["class"] == "spam").astype(int)
X_tr, X_te, y_tr, y_te = train_test_split(X, y, test_size=0.3,
                                          stratify=y, random_state=0)

# vetorizacao + classificador num unico objeto (sem vazamento!)
pipe = Pipeline([
    ("vec", TfidfVectorizer(lowercase=True, stop_words="english")),
    ("clf", MultinomialNB()),
])
grade = {"vec__ngram_range": [(1,1), (1,2)],
         "clf__alpha":       [0.1, 0.5, 1.0]}     # suavizacao de Laplace
busca = GridSearchCV(pipe, grade, cv=5, scoring="f1").fit(X_tr, y_tr)

y_hat = busca.predict(X_te)
print(confusion_matrix(y_te, y_hat))
print(classification_report(y_te, y_hat))        # precisao, recall, F1
\end{lstlisting}
O notebook \texttt{Aula prática E3} percorre a limpeza, a exploração
(palavras mais comuns por classe) e a classificação na base \texttt{spam.csv}.

%==============================================================================
\section*{Resumo (e encerramento do curso)}
%==============================================================================
\begin{itemize}[leftmargin=1.4cm]
  \item Texto vira número por \textbf{\emph{bag-of-words}} / \textbf{TF--IDF} após
        limpeza (minúsculas, pontuação, \emph{stopwords}).
  \item A matriz documento--termo é \textbf{esparsa e de alta dimensão} ---
        terreno do \textbf{Bayes ingênuo} e da \textbf{logística penalizada}
        (Aulas~05 e~08).
  \item Avalie com \textbf{precisão/\emph{recall}/$F_1$} e ajuste o \textbf{corte}
        (dados desbalanceados, Aula~09).
  \item Tudo num \textbf{\emph{pipeline}} com CV sem vazamento (Aulas~03 e~07).
\end{itemize}

\begin{ideia}[O fio condutor do curso]
Este exemplo mostra que os ``temas'' não são ilhas: representar dados, controlar o
balanço viés--variância, estimar o risco honestamente, escolher o método pela
estrutura do problema e medir o desempenho com a métrica certa. Foi essa a história
do curso inteiro --- da regressão linear ao \emph{spam}.
\end{ideia}

\paragraph{Leitura recomendada.} [AME] §8.1.3 (Bayes ingênuo), §3.8 e §8.8 (exemplos
com texto: resenhas da Amazon), Apêndice A.2 (manipulação de textos) e A.3 (matrizes
esparsas); documentação do \texttt{scikit-learn}: \emph{Working With Text Data}.

\end{document}
